\chapter{Ureteral Stent Placement}

\section*{Overview}
Ureteral stent placement is the insertion of a thin, flexible tube into the ureter to maintain drainage of urine from the kidney to the bladder. It is placed when a ureter is obstructed or when drainage is needed after urologic procedures.
\section*{Alternative Names}
Ureteral stent insertion; double-J stent; JJ stent; retrograde ureteral stent
\section*{Principle}
A stent with curled ends (double-J) is passed through the ureter so that its coils rest in the renal pelvis and bladder, holding the ureter open. Urine then drains by gravity and peristalsis around and through the stent.
\section*{History}
Temporary ureteral drainage was refined through the 20th century, and the double-J stent introduced by Finney in 1978 became the standard device for maintaining ureteral patency.
\section*{Why It Is Done}
Ordered for ureteral obstruction from stones, tumors, or strictures, to protect the kidney during and after ureteroscopy or lithotripsy, to bypass an injured ureter, and to relieve hydronephrosis when urgent drainage is required.
\section*{Is This a Primary or Secondary Test or Procedure}
Primary therapeutic procedure for ureteral obstruction and a supportive measure during other urologic procedures.
\section*{How It Is Performed}
Under fluoroscopic guidance, a guidewire is passed through a cystoscope into the ureter, and the stent is advanced over the wire until its proximal coil is in the renal pelvis and its distal coil in the bladder. The wire is removed and correct position is confirmed by imaging.
\section*{Preparation}
Informed consent, review of anticoagulants and bleeding risk, and evaluation of renal function and urine culture before the procedure.
\section*{Contraindications and Precautions}
Contraindications include untreated urinary tract infection and severe bleeding diathesis. Precautions include reviewing the patient for contrast allergy if imaging with contrast is used and using caution in pregnancy.
\section*{Risks}
Pain and hematuria, bladder irritation and urinary frequency, stent migration or obstruction, infection, and rarely ureteral injury or perforation.
\section*{Aftercare and Recovery}
The patient is monitored for hematuria and voiding symptoms and given analgesia as needed. The stent is left in place for the planned duration and removed endoscopically.
\section*{Outcome and Follow-up}
Good urine drainage with resolution of hydronephrosis and relief of obstruction indicates success. Persistent pain, fever, or poor output may indicate stent dysfunction or infection.
\section*{Expected Results}
A correctly positioned stent with free drainage of urine, no infection, and resolution of obstruction.
\section*{Follow-up}
The stent is removed or exchanged at the scheduled interval, typically within weeks to months, and imaging is repeated to confirm drainage of the kidney.
\section*{References}
\begin{enumerate}
  \item MedlinePlus. Ureteral stent. U.S. National Library of Medicine; 2025.
  \item Current Medical Diagnosis and Treatment. McGraw-Hill; 2025.
\end{enumerate}

\clearpage
